\documentclass[11pt]{article}

\usepackage{mathtools}

\DeclareMathOperator\vectorize{vec}


\title{Rigidity Matrices}
\author{W. Muneo Okiishi}
\date{2026-10-15}

\begin{document}

\maketitle

\section{Notation}

We denote matrices with bold capital letters, such as $\mathbf{M}$, and
vectors with bold lowercase letters, such as $\mathbf{v}$. We denote the
elements thereof with subscripts and without boldface, for instance $M_{ij}$
or $v_k$.

If $\mathbf{M}$ is a matrix, we use $\mathbf{M}_{i*}$ to denote its $i$-th row
vector and $\mathbf{M}_{*j}$ to refer to its $j$-th column vector.
$\mathbf{M}_{i*}^\mathrm{T}$ and $\mathbf{M}_{*j}^\mathrm{T}$ denote the
transposes of these respective vectors.

We use $\otimes$ to denote the Kronecker product.

We use $\left \lVert \mathbf{v} \right \rVert$ to denote the standard Euclidean
2-norm of a vector $\mathbf{v}$.

We use $\vectorize(\mathbf{M})$ to refer to the column vector produced by a
row-major vectorization of a matrix $\mathbf{M}$.

\section{Constructing a Rigidity Matrix}

Let $\mathbf{B}$ be the directed incidence matrix for a framework,
defined with rows corresponding to links and columns corresponding to nodes.
(Note that this is the transpose of the more common definition.)

Let $\mathbf{X}$ be the position matrix of the same framework.
It comprises a stack of row vectors, each corresponding to the position of a
node in the framework.

We can now define an intermediate matrix $\mathbf{V}$, representing the
direction of each link in the framework:
\begin{equation} \label{eq:V_def}
    \mathbf{V} = \mathbf{B} \mathbf{X}
\end{equation}

Using this, each row of the rigidity matrix $\mathbf{R}$ of the framework can
be expressed as:
\begin{equation} \label{eq:R_def}
    \mathbf{R}_{i*} = \mathbf{B}_{i*} \otimes \mathbf{V}_{i*}
\end{equation}

Algorithms constructing $\mathbf{R}$ can take advantage of a variety of
identities. For example:
\begin{align}
    \mathbf{R}_{i*}^\mathrm{T}
    &= \vectorize(\mathbf{B}_{i*}^\mathrm{T} \mathbf{V}_{i*}) \\
    &= \vectorize(\mathbf{B}_{i*}^\mathrm{T} \mathbf{B}_{i*} \mathbf{X}) \\
    &= ((\mathbf{B}_{i*}^\mathrm{T} \mathbf{B}_{i*}) \otimes \mathbf{I})
    \vectorize(\mathbf{X}) \label{eq:R_expanded}
\end{align}

\section{Normalized Rigidity Matrix}

Some applications require a normalized rigidity matrix $\mathbf{\hat{R}}$.
To calculate this, we first construct a new matrix $\mathbf{\hat{V}}$ by
normalizing each row of $\mathbf{V}$:
\begin{equation} \label{eq:Vn_def}
    \mathbf{\hat{V}}_{i*}
    = \frac{\mathbf{V}_{i*}}{\left \lVert \mathbf{V}_{i*} \right \rVert}
\end{equation}

With this, calculating the normalized rigidity matrix takes the same form as \eqref{eq:R_def}:
\begin{equation} \label{eq:Rn_def}
    \mathbf{\hat{R}}_{i*} = \mathbf{B}_{i*} \otimes \mathbf{\hat{V}}_{i*}
\end{equation}

$\mathbf{\hat{R}}$ can thus be constructed with similar algorithms as $\mathbf{R}$.

\section{Differentiating the Normalized Rigidity Matrix}

Some applications may require $\frac{\partial \mathbf{\hat{R}}}{\partial x_k}$.
We can derive this by differentiating each row of $\mathbf{\hat{R}}$.

From \eqref{eq:Vn_def} and \eqref{eq:Rn_def}, we can see that
\begin{equation}
    \mathbf{\hat{R}}_{i*} = a_i \mathbf{R}_{i*}
\end{equation}
where
\begin{equation} \label{eq:a_def}
    a_i = \left \lVert \mathbf{V}_{i*} \right \rVert^{-1}
    = (\mathbf{V}_{i*} \mathbf{V}_{i*}^\mathrm{T})^{-1/2}
\end{equation}

This means that
\begin{equation} \label{eq:dRndx_basic}
    \frac{\partial \mathbf{\hat{R}}_{i*}}{\partial x_k}
    = a_i \frac{\partial \mathbf{R}_{i*}}{\partial x_k}
    + \frac{\partial a_i}{\partial x_k} \mathbf{R}_{i*}
\end{equation}

From \eqref{eq:R_expanded}, we can derive $\frac{\partial \mathbf{R}_{i*}}{\partial x_k}$:
\begin{equation} \label{eq:dRdx}
    \frac{\partial \mathbf{R}_{i*}}{\partial x_k}
    = ((\mathbf{B}_{i*}^\mathrm{T} \mathbf{B}_{i*}) \otimes \mathbf{I})_{*k}^\mathrm{T}
\end{equation}

To find $\frac{\partial a_i}{\partial x_k}$, we start by differentiating
\eqref{eq:a_def} with respect to $X_{nd}$:
\begin{equation}
    \frac{\partial a_i}{\partial X_{nd}}
    = -a_i^3 \frac{\partial \mathbf{V}_{i*}}{\partial X_{nd}} \mathbf{V}_{i*}^\mathrm{T}
\end{equation}

Referencing \eqref{eq:V_def}, we can convert this to
\begin{equation}
    \frac{\partial a_i}{\partial X_{nd}}
    = -a_i^3 \mathbf{B}_i \frac{\partial \mathbf{X}}{\partial X_{nd}} \mathbf{V}_{i*}^\mathrm{T}
\end{equation}

This reduces to
\begin{equation}
    \frac{\partial a_i}{\partial X_{nd}} = -a_i^3 B_{in} V_{id}
\end{equation}

Switching from $X_{nd}$ to $x_k$, this simplifies to
\begin{equation} \label{eq:dadx_simple}
    \frac{\partial a_i}{\partial x_k} = -a_i^3 R_{ik}
\end{equation}

Finally, we can substitute \eqref{eq:dadx_simple} into \eqref{eq:dRndx_basic}:
\begin{equation}
    \frac{\partial \mathbf{\hat{R}}_{i*}}{\partial x_k}
    = a_i \frac{\partial \mathbf{R}_{i*}}{\partial x_k} - a_i^3 R_{ik} \mathbf{R}_{i*}
\end{equation}

An algorithm implementing this to construct $\frac{\partial \mathbf{\hat{R}}}{\partial x_k}$
can be optimized in a variety of ways depending on the application.

\end{document}
